\section{AI Model Deployment}

\subsection{Model Registry}

Model Registry là kho quản lý metadata và artifact của các phiên bản mô hình.
Mỗi bản ghi liên kết model name, version, framework, precision, input shape, chỉ
số đánh giá, vị trí artifact, checksum và ma trận tương thích thiết bị. Registry
không chỉ là thư mục lưu tệp mà là nguồn sự thật cho quyết định triển khai.

\subsection{Model Versioning}

Model versioning giúp phân biệt các checkpoint, cấu hình tối ưu và engine. Một
phiên bản đã phát hành phải bất biến; thay đổi nội dung yêu cầu tạo phiên bản mới.
Ngoài version dễ đọc, hệ thống lưu checksum để nhận diện chính xác bytes của
artifact và ngăn trường hợp hai tệp khác nhau dùng chung tên.

\subsection{Model Artifact}

Artifact triển khai có thể gồm TensorRT engine, ONNX dự phòng, label map,
manifest, cấu hình preprocessing--postprocessing và chữ ký hoặc checksum. Gói
phải tự mô tả đủ để Edge Agent kiểm tra tính tương thích trước khi dừng runtime
hiện tại.

\subsection{Deployment Strategy}

Chiến lược triển khai xác định thứ tự và phạm vi node nhận model. Các lựa chọn
gồm triển khai đồng loạt, rolling deployment và canary. Với đội thiết bị phân tán,
canary giúp kiểm tra một nhóm nhỏ trước khi mở rộng, đồng thời đặt ngưỡng dừng
dựa trên health check, FPS, error rate hoặc nhiệt độ.

\subsection{Rollback}

Rollback đưa node về phiên bản đã biết là ổn định khi download, validation,
activation hoặc health check thất bại. Edge Agent cần giữ ít nhất một artifact
trước đó, cập nhật con trỏ active theo cách nguyên tử và báo cáo nguyên nhân. Cơ
chế rollback phải được kiểm thử như một luồng bình thường, không chỉ dùng khi
sự cố đã xảy ra.